\documentclass[10pt]{article} 
\renewcommand{\thempfootnote}{\arabic{mpfootnote}}

\usepackage[utf8]{inputenc}

\usepackage{multicol}
\setlength{\columnsep}{1cm}

\usepackage{abc}
\usepackage{graphicx}
\usepackage{musicography}

\usepackage[margin=0.5in]{geometry} 
\geometry{letterpaper}

\usepackage{fancyhdr} 
\pagestyle{fancy} % options: empty , plain , fancy
\renewcommand{\headrulewidth}{0pt} % customise the layout...
\lhead{}\chead{}\rhead{}
\lfoot{}\cfoot{\thepage}\rfoot{}

\begin{document}

\thispagestyle{empty}

\begin{abc}[name=sallyann,options={-O= +c -F single.fmt}]
X:1
T:Sally Ann
K:A
M:C|
|: ef | "A"a2 ab af e2 | "D"f2-ff f4 | a2 f2 ec-c2 | "A"e2-ee e2 ef | 
e2 c2 B2 A2 | "E"B2-BB B4 | cBAc BAFE | "A"FA-AB A2 :|
|: AF | "A"EFAB cBAc | "E"B2 BA B4 | cBAc BAFE | "A"FA-AB A4 | 
EFAB cBAc | "E"B2 BA B2 cd | efed cABc | "A"A2 AB A2 :|
\end{abc}
% \vskip 0.5in
% \noindent\makebox[\textwidth][c]{%
%   \begin{minipage}{7.0in}
%    \setlength{\parindent}{1.5em}
    
    The tune, Sally Ann, stems from the same tune family as Sail Away
    Ladies and Great Big Taters in Sandy Land, sounding like variations
    of each other with the main difference being what area of the
    country they were predominantly played. All three tunes are thought
    to be from the mid-19th century, but which tune came first and from
    where is unclear.\footnote{https://tunearch.org/wiki/Annotation:Great\_Big\_Taters\_in\_Sandy\_Land}
    
    Sally Ann is associated with West Virginia, Virginia, and North
    Carolina; being so popular in Surry County, NC, that it is
    considered the unofficial `national
    anthem'.\footnote{https://tunearch.org/wiki/Annotation:Sally\_Ann\_(1)}
    It is played both in standard tuning (GDAE) and crosstuning (ADAE)
    with famous instrumental versions performed by Doc Roberts in
    1931\footnote{https://youtu.be/bVMH45dqSWI} and later, Lester Flatt
    and Earl Scruggs in 1961.\footnote{https://youtu.be/zpv\_8xEJFYg}
    
    The song, Sally Ann, makes appearances in English Folk Songs from
    the Southern Appalachians Vol 2, collected by Cecil Sharp and edited
    by Maud Karpeles.\footnote{Sharp, Cecil J. and Maud
      Karpeles. English Folk Songs from the Southern Appalachians Vol 2,
      Oxford University Press, 1932, No 240, p 351.} The tune was sung
    to Sharp and Karpeles by Mrs.~Dellie Hughes on October 9th, 1918 in
    Burnsville, NC.
    
    \begin{quote}
      O where are you going, Sally Anne? (x3)\\
      I'm going to the wedding, Sally Anne.
      \\
      O shake that little foot, Sally Anne, (x3)\\
      You're a pretty good dancer, Sally Anne.
    \end{quote}
    
    It is also included in Folk Songs of North America, Alan
    Lomax,\footnote{Lomax, Alan. Folk Songs of North America, Doubleday
      and Company, 1960, p 230.} published in 1960, but it is not certain
    where or when it was gathered.
    
    \begin{multicols}{2}
      \noindent
      Did you ever see a muskrat, Sally Anne, \\
      Draggin' his slick tail through the sand,\\
      Pickin' his banjo an' raisin' sand,\\
      Did you ever see a muskrat, Sally Anne?
      \vfill\null
      \columnbreak
      \noindent
      Chorus: \\
      Ever see a muskrat, Sally, Sally\\
      Ever see a muskrat, Sally Anne?
    \end{multicols}
    \vspace*{-.5cm}
    Sally Ann makes a more recent appearance in 1990, when the New
    Nashville Cats, fronted by John Cowan and Mark O'Connor, performed a
    new arrangement of Sally Ann with lyrics written by Alice Randall
    that include a race lynching.\footnote{https://youtu.be/B4-52uUYoK8}
    
    \begin{multicols}{2}
      \noindent
      Johnny got married in his one good suit \\
      But the ride from church bore strange fruit\\
      Down by road you can hear her cry \\
      As he hung from a tree, she watched him die
      \vfill\null
      \columnbreak
      \noindent
      Refrain:\\
      Who's gonna dance with Sally Ann \\
      Who's gonna touch her tremblin' hand \\
      When the fiddler takes the stand \\
      Who's gonna dance with Sally Ann
    \end{multicols}
%  \end{minipage}
% }


\end{document}
